% !TeX encoding = UTF-8
\chapter{Discussion and Conclusion}\label{chap:discussion}

\paragraph{Shortcut features and domain shift.} The near-perfect benchmark scores (DDoS F1 0.9996) may partly rest on a shortcut. \texttt{Destination Port} ranks second among the selected features, and virtually all DDoS flows target port 80, so the model may have learned the port rather than the attack behavior (Section~\ref{sec:t3-interp}). The test split comes from the same capture as the training data, so it cannot reveal such a shortcut. The CICFlowMeter version, the hosts and the application mix also differ between the 2017 testbed and our network (Section~\ref{sec:latency}). The 0 alerts on our capture show only that this particular benign traffic is not mistaken for DDoS. The data itself also has quality problems, such as the duplicated column, infinite values, 2,633 duplicate rows and several TCP flag columns that are identical or constant. They match issues reported in the literature~\cite{engelen2021troubleshooting}, and each of them needed its own cleaning step in Chapter~\ref{chap:task3_preprocessing}.

\paragraph{Class balance.} Because the classes are close to balanced, class weighting has almost no effect: the balanced and default Decision Trees are identical on the test split, and the balanced Logistic Regression removes only 34 of about 2,000 false alarms.

\paragraph{Limitations.} We use a single train/test split without cross-validation. The tree is unpruned. We ran no ablation without \texttt{Destination Port}, so we cannot quantify the shortcut. Dropping rows with infinite values removes 34 flows, including 2 DDoS flows, and min-max scaling is sensitive to outliers. Finally, the real-time part replays recorded flows and is not live capture.

\paragraph{Possible improvements (not implemented).} Pruning the tree (\texttt{max\_depth}, \texttt{min\_samples\_leaf}), dropping the testbed-specific features (such as \texttt{Destination Port}), evaluating on a second attack capture from another network, and adding benign traffic from the deployment network to the training data would make the conclusions more reliable.

\paragraph{Conclusion.} We turned the Friday DDoS file of CIC-IDS2017 into a clean set of 223,082 flows with 59 features, selected the 20 best features, and trained Logistic Regression and Decision Tree classifiers for BENIGN vs.\ DDoS. The Decision Tree outperforms Logistic Regression and makes only 32 errors on 66,925 test flows. On our own benign traffic, it raised no false alarm, but because our capture has no attack and the data may allow shortcuts, its benchmark score does not by itself show how well it would detect DDoS attacks in our network.
